% BEGIN CR28_RADIAL_FULL
% Sección incorporable; requiere amsmath y amssymb.
% Fuente principal: RIGIDEZ_SIMULTANEA_DEL_LECTOR_RADIAL_20260926.md.
% Complementos demostrativos: COMPOSICION_METRICA_MEMORIA_Y_ACOPLAMIENTO.md
% y ACCION_ACOPLADA_FASE_Y_MEMORIA_DE_FRONTERA.md.
% La integración, la compilación y la revisión visual corresponden al editor.
\section{Rigidité simultanée du lecteur longitudinal et du couplage radial}
\label{g26:sec:rigidez-radial-es}

La construction de l'holographie modulaire triadique transmet à la réalisation
radiale un état enrichi avec ses coordonnées, son orientation, sa retenue,
sa mémoire et ses incidences. La suite
\[
 \begin{gathered}
 \mathrm{APP}\longrightarrow\mathrm{TRIT}\longrightarrow\mathrm{TPK}
 \longrightarrow\text{état enrichi}\\
 \longrightarrow\text{structure discrète conjointe du continu}
 \end{gathered}
\]
conserve la projection archimédienne et la projection d'incidence, ainsi que les cinq constructions consubstantielles du continu. Le prolongement
\[
 w_6\longrightarrow w_{12}\longrightarrow w_{18}\longrightarrow
 w_{24}\longrightarrow w_{30}\longrightarrow R_{36}\longrightarrow G_9
\]
s'articule avec l'ordre opératoriel
\(U_t\longrightarrow\Gamma_9\longrightarrow K_{\rm ph}\) : la phase revient
et le registre de mémoire s'incrémente. Les coordonnées régionales, le
registre sélectionné \(K\), la coordonnée \(\alpha\) et l'action de retour
sont reçus avec cette provenance commune et leurs lecteurs préalablement
construits.

Le résultat de cette section est une classification de réalisations qui conservent simultanément les données générées et les règles longitudinales, métriques et temporelles précisées ci-après. Pour chaque caractère positif, la réponse radiale est déterminée ; le coefficient qui relie cette réponse à la masse est commun à tous les caractères admissibles. L'interprétation gravitationnelle est formulée après la construction des deux lectures sur le même espace.

\subsection{Registre des incidences et contraintes constitutives}

On considère des espaces de Hilbert complexes \(V_{\rm in}\) et
\(V_{\rm out}\) de dimension \(N\), avec leurs produits scalaires fixés.
Le cardinal marqué spécifié dans R1 est positif ; en particulier,
les deux fibres sont non nulles. Les adjoints sont pris par rapport à ces produits. Le symbole
\(\mathsf U\) désignera le transport normalisé de l'archive et
\(U_t\) conservera sa signification d'évolution de l'état enrichi.

\paragraph{R1. Registre marqué.}
L'ensemble des incidences est
\begin{equation}
 \Omega=\left\{(j,k,q,u,v):
 \begin{array}{l}
 j,k\in\mathbb Z/12\mathbb Z,\quad k-j\in\{0,3,6,9\},\\
 1\leq q\leq8,\quad u<v,\quad
 u,v\in\{1,2,4,5,7,8\}
 \end{array}\right\}.
 \label{g26:eq:omega-rigidez-es}
\end{equation}
Les douze phases, les quatre positions relatives, les huit positions de l'
indice \(q\) et les quinze paires du dernier ensemble donnent
\begin{equation}
 N=|\Omega|=12\cdot4\cdot8\binom62
   =12\cdot4\cdot120=5760.
 \label{g26:eq:cardinal-rigidez-es}
\end{equation}
Chaque incidence possède sa marque individuelle. Dans la base correspondante \((e_i)_{i\in\Omega}\) de \(V_{\rm out}\), soient \(P_i=e_ie_i^*\). Le mode commun \(\boldsymbol u\) est normalisé par \(|u_i|^2=1/N\) ; il s'écrit \(P=\boldsymbol u\boldsymbol u^*\). L'algèbre marquée est \(\mathcal A_\Omega=\operatorname{span}\{P_i:i\in\Omega\}\).

\paragraph{R2. Normalisation et archive complète.}
L'inverseur généré \(B:V_{\rm in}\to V_{\rm out}\) satisfait
\begin{equation}
 B^*B=\tfrac14 I_{\rm in},\qquad
 BB^*=\tfrac14 I_{\rm out}.
 \label{g26:eq:normalizacion-b-es}
\end{equation}
Les deux composantes sont conservées
\begin{equation}
 C_0=(I-P)B,\qquad M_0=PB,\qquad
 C_0+M_0=B,\qquad \mathsf U=2B.
 \label{g26:eq:archivo-completo-es}
\end{equation}
La première enregistre la variation par rapport au mode commun ; la seconde
conserve ce mode. Leurs images sont orthogonales et leur somme restitue
le transport complet. D'après \eqref{g26:eq:normalizacion-b-es},
\(\mathsf U\) est unitaire.

\paragraph{R3. Inscription individuelle.}
L'application \(\mathcal E:\operatorname{End}(V_{\rm out})\to\mathcal A_\Omega\) est une espérance conditionnelle linéaire qui fixe chaque \(P_i\) et satisfait
\begin{equation}
 \mathcal E(A_1TA_2)=A_1\mathcal E(T)A_2,
 \qquad A_1,A_2\in\mathcal A_\Omega.
 \label{g26:eq:bimodularidad-es}
\end{equation}
L'inscription du résultat conserve, dans un registre complémentaire,
les données complètes de la construction. L'espérance détermine
l'observable inscrit ; la conservation de l'archive détermine sa
reconstruction ultérieure.

\paragraph{R4. Composition longitudinale.}
On fixe la coordonnée générée \(\alpha>0\) et la section de référence longitudinale \(L_*>0\). Le transport longitudinal est défini par
\begin{equation}
 D=L_*\alpha^{16}\mathcal E(C_0C_0^*)\mathsf U:
 V_{\rm in}\longrightarrow V_{\rm out}.
 \label{g26:eq:regla-longitudinal-es}
\end{equation}
Cette règle applique linéairement le coefficient inscrit à la section
longitudinale. Son type fixe la position de la puissance \(\alpha^{16}\)
et du coefficient de retour dans la composition. La lecture d'
une racine d'intensité constitue une opération différente et possède sa
propre règle de transformation.

\paragraph{R5. Métrique radiale positive.}
Soit \(X=X^*>0\) le caractère complet sur \(V_{\rm in}\). On définit
la lecture radiale positive \(R_X\) sur \(V_{\rm out}\) par
\begin{equation}
 \|R_Xv\|^2=\|XD^*v\|^2
 \qquad(v\in V_{\rm out}).
 \label{g26:eq:metrica-radial-es}
\end{equation}
Cette égalité fixe la métrique sur chaque vecteur de la fibre. En notation polaire, \(R_X\) est le module d'arrivée \(\lvert(DX)^*\rvert=\sqrt{DX^2D^*}\).

\paragraph{R6. Action et horloge.}
La section d'action générée \(\hbar_{\rm ret}>0\) et le relèvement
réel de phase satisfont
\begin{equation}
 S(\theta)=\hbar_{\rm ret}\theta.
 \label{g26:eq:accion-fase-es}
\end{equation}
La vitesse constitutive \(c>0\) est conservée. Une fois obtenue la longueur intensive \(L\) de \eqref{g26:eq:regla-longitudinal-es}, l'horloge conjuguée satisfait \(\omega L=c\). Pour \(\theta(t)=\omega t\), le taux d'action détermine l'échelle énergétique
\begin{equation}
 E_{L}=\frac{dS}{dt}=\hbar_{\rm ret}\omega
   =\frac{\hbar_{\rm ret}c}{L}.
 \label{g26:eq:escala-energetica-es}
\end{equation}
Le relèvement et ses raffinements restent associés à la mémoire
de la route.

\paragraph{R7. Lectures du même caractère.}
Sur \(Y=\mathsf U X\mathsf U^*>0\), on définit
\begin{equation}
 H_X=E_{L}Y,\qquad M_X=\frac{H_X}{c^2},\qquad
 \Lambda_X=LY^{-1}.
 \label{g26:eq:lecturas-conjugadas-es}
\end{equation}
L'énergie, la masse et la longueur conjuguée conservent le même caractère
et le même transport. L'inverse agit dans le secteur positif ; le secteur
nul conserve son traitement propre.

\paragraph{R8. Rayon réduit et sections dimensionnelles.}
L'identification physique de cette réalisation est
\begin{equation}
 R_X=\frac{G}{c^2}M_X.
 \label{g26:eq:radio-reducido-es}
\end{equation}
Ainsi, \(R_X\) représente le rayon gravitationnel réduit. Le rayon de Schwarzschild correspondant est \(2R_X\). On maintient conjointement \(L_*\), la section d'action et la section constitutive de vitesse ; un changement d'unités transforme leurs coordonnées selon leurs dimensions.

\paragraph{Provenance des restrictions.}
Le transport de Hadamard, la conservation de l'archive, l'inscription,
l'action de phase et la réciprocité possèdent des antécédents dans le corpus.
Le registre élargi \eqref{g26:eq:omega-rigidez-es} et les compositions
\eqref{g26:eq:regla-longitudinal-es}--\eqref{g26:eq:metrica-radial-es} sont
incorporés explicitement dans la construction réunie. R4 et R5 sont des règles
constitutives de cette réalisation ; la démonstration suivante établit
leur conséquence conjointe avec les autres antécédents. R8 précise
l'identification physique qui fixe le sens du coefficient.

\subsection{Théorème de détermination simultanée}

\paragraph{Théorème.}
Pour les mêmes données générées \(B\), \((P_i)\), \(P\),
\(\alpha\), \(L_*\), \(\hbar_{\rm ret}\), \(c\) et le même caractère
\(X\), les contraintes R1--R8 déterminent univoquement l'inscription,
le transport longitudinal et la réponse radiale. Pour tous les
caractères positifs admissibles, le coefficient \(G\) est commun. On a
\begin{align}
 \mathcal E&=\Delta_\Omega,
 &\Delta_\Omega(T)&=\sum_{i\in\Omega}P_iTP_i,\label{g26:eq:delta-unica-es}\\
 r_N&=\frac{N-1}{4N}=\frac{5759}{23040},
 &L&=L_*\alpha^{16}r_N,\label{g26:eq:longitud-unica-es}\\
 D&=L\mathsf U,
 &R_X&=L\mathsf U X\mathsf U^*,\label{g26:eq:radial-unica-es}\\
 R_X\Lambda_X&=L^2I,
 &H_X\Lambda_X&=\hbar_{\rm ret}cI.\label{g26:eq:productos-radiales-es}
\end{align}
Le couplage est déterminé par
\begin{equation}
 \boxed{G=\frac{c^3L^2}{\hbar_{\rm ret}}
 =\frac{c^3L_*^2}{\hbar_{\rm ret}}
   \left(\frac{5759}{23040}\right)^2\alpha^{32}.}
 \label{g26:eq:g-rigido-es}
\end{equation}

\paragraph{Démonstration : unicité de l'inscription.}
Soient \(E_{ij}=e_ie_j^*\) les unités matricielles. Par bimodularité,
\[
 \mathcal E(E_{ij})=P_i\mathcal E(E_{ij})P_j.
\]
Pour \(i\ne j\), l'appartenance de \(\mathcal E(E_{ij})\) à l'algèbre diagonale implique \(P_i\mathcal E(E_{ij})P_j=0\). Pour \(i=j\), la fixation de \(P_i\) donne \(\mathcal E(E_{ii})=P_i\). La linéarité détermine alors \eqref{g26:eq:delta-unica-es} sur toutes les matrices. L'inscription marquée est ainsi fixée par ses propriétés algébriques.

\paragraph{Démonstration : coefficient de retour et archive.}
La normalisation de \(B\) et \(P^2=P=P^*\) fournissent
\[
 C_0C_0^*=\tfrac14(I-P),\qquad M_0M_0^*=\tfrac14P.
\]
Comme \(P_iPP_i=|u_i|^2P_i=P_i/N\), on obtient
\begin{equation}
 \Delta_\Omega(C_0C_0^*)=\frac{N-1}{4N}I,
 \qquad
 \Delta_\Omega(M_0M_0^*)=\frac{1}{4N}I.
 \label{g26:eq:retorno-y-memoria-es}
\end{equation}
Les deux lectures ont pour somme \(I/4\). La première est scalaire, de sorte que
toute distribution normalisée \((p_i)\) sur les incidences,
avec \(p_i\geq0\) et \(\sum_i p_i=1\), produit
\(\sum_i p_i r_N=r_N\). Le même coefficient est donc préservé
sous une pondération diagonale arbitraire.

\paragraph{Démonstration : longueur et métrique.}
La substitution de \eqref{g26:eq:retorno-y-memoria-es} dans R4 donne
\(D=L\mathsf U\), avec \(L\) comme dans \eqref{g26:eq:longitud-unica-es}.
L'unitarité de \(\mathsf U\) implique \(DD^*=L^2I\).
D'après R5, les formes quadratiques de \(R_X^2\) et \(DX^2D^*\) coïncident.
L'identité de polarisation détermine
\[
 R_X^2=DX^2D^*=L^2\mathsf U X^2\mathsf U^*.
\]
L'opérateur \(L\mathsf U X\mathsf U^*\) est positif et son carré est l'opérateur du second membre. L'unicité de la racine carrée positive démontre \eqref{g26:eq:radial-unica-es}. En particulier, \(R_I=LI\).

\paragraph{Démonstration : action, réciprocité et couplage.}
R6 et R7 donnent
\[
 H_X=\frac{\hbar_{\rm ret}c}{L}Y,
 \qquad M_X=\frac{\hbar_{\rm ret}}{cL}Y,
 \qquad \Lambda_X=LY^{-1}.
\]
La multiplication produit les deux produits de
\eqref{g26:eq:productos-radiales-es}. Ensuite, R8 équivaut à
\[
 LY=\frac{G\hbar_{\rm ret}}{c^3L}Y.
\]
L'inversibilité de \(Y\) détermine \eqref{g26:eq:g-rigido-es}.
La même substitution vérifie l'égalité pour chaque caractère positif.
Si \(G_1\) et \(G_2\) satisfont l'identité avec les mêmes données,
\((G_1-G_2)Y=0\), et donc \(G_1=G_2\).
Cette conclusion distingue la réponse dépendant de l'état
\(R_X\) du coefficient commun \(G\). \hfill\(\square\)

\subsection{Réciprocité et rigidité de la phase}

L'antécédent puissance--logarithmique utilise, pour \(x>0\),
\[
 \ell_p(x)=\frac{x^p-1}{p}\quad(p\ne0),
 \qquad \ell_0(x)=\log x.
\]
La substitution directe démontre
\(\ell_{-p}(x^{-1})=-\ell_p(x)\). Dans le domaine positif de leurs
inverses, cette identité équivaut à
\[
 \exp_p(s)\exp_{-p}(-s)=1.
\]
La composition spectrale sur un caractère positif transmet ce produit à la paire \(LY\), \(LY^{-1}\). Dans la réalisation actuelle, R5 peut s'exprimer de manière équivalente, une fois R1--R4 et R7 fixées, par \(R_X\Lambda_X=L^2I\) : multiplier par \(\Lambda_X^{-1}\) détermine \(R_X=L\mathsf U X\mathsf U^*\), et cette expression satisfait R5.

Les deux produits de \eqref{g26:eq:productos-radiales-es} donnent en outre une
forme directe de la rigidité :
\begin{equation}
 R_X=\frac{L^2}{\hbar_{\rm ret}c}H_X.
 \label{g26:eq:proporcion-operadores-es}
\end{equation}
Pour un changement d'échelle positif \(s\), le couple
\((sR_X,\Lambda_X/s)\) conserve le produit radial ; le produit
\(H_X(\Lambda_X/s)=\hbar_{\rm ret}cI/s\), à énergie et action
fixées, exige \(s=1\). La conservation conjointe détermine l'échelle.
De même, tout opérateur radial positif ayant la même métrique R5
coïncide avec \(R_X\) par l'unicité de la racine positive.

L'action relevée possède aussi une détermination exacte. Pour
\(\theta\ne0\), l'égalité \(S/\hbar_{\rm ret}=\theta\) fixe \(S\).
Même lorsque les rotations de toutes les subdivisions
\(\theta/9^n\) sont conservées, un changement d'échelle réel fixe \(a\) qui satisfait
\[
 \exp(-ia\theta/9^n)=\exp(-i\theta/9^n)
 \qquad(n\geq0)
\]
vérifie \((a-1)\theta/9^n\in2\pi\mathbb Z\). Cette suite tend vers zéro ; pour \(n\) suffisamment grand, sa seule valeur possible est zéro. De \(\theta\ne0\), on déduit \(a=1\).

\subsection{Échelles de Planck et spécialisation du caractère unitaire}

Nous écrivons \(L_\alpha=L\) ; dans la section de retour, \(E_{P,\rm ret}=E_L\).

La longueur construite $L_\alpha>0$, la section d'action $\hbar_a>0$ et la vitesse $c>0$ déterminent les échelles

\[
t_P=\frac{L_\alpha}{c},\qquad
m_{P,a}=\frac{\hbar_a}{cL_\alpha},\qquad
E_{P,a}=\frac{\hbar_a c}{L_\alpha}.
\]

En effet, substituer $G_a=c^3L_\alpha^2/\hbar_a$ dans les expressions
$\sqrt{\hbar_aG_a/c^5}$, $\sqrt{\hbar_ac/G_a}$ et
$\sqrt{\hbar_ac^5/G_a}$ donne respectivement ces trois identités.
La positivité fixe dans chaque cas la racine. La longueur précède la
reconnaissance de ces échelles ; leurs expressions conventionnelles sont
des identités de la réalisation déjà construite.

Pour un mode de caractère $y>0$, les lectures massique et longitudinale sont

\[
m(y)=m_{P,a}y,\qquad r_g(y)=L_\alpha y,\qquad
\bar\lambda(y)=\frac{L_\alpha}{y}.
\]

Par conséquent, $r_g(y)\bar\lambda(y)=L_\alpha^2$ pour tous les modes
positifs. L'égalité des deux longueurs caractérise le mode $y=1$,
puisque $y=y^{-1}$ et $y>0$ impliquent $y=1$. Dans ce mode, la masse est
$m_{P,a}$ et les deux longueurs coïncident avec $L_\alpha$. On distingue ainsi
la réciprocité générale de sa spécialisation au caractère unitaire.

Le pas chronologique et le temps de Planck satisfont $t_0=(\pi_{\rm HMT}/54)t_P$. La période complète de phase de 108 pas est $108t_0=2\pi_{\rm HMT}t_P$ et sa fréquence angulaire est $\omega_{108}=1/t_P=c/L_\alpha$. Un horizon de rayon $R_h=\zeta_hL_\alpha$, avec $\zeta_h>0$, conserve sa propre variable géométrique. Le rayon de Schwarzschild associé à une masse est $2r_g$.

Si l'on compare les deux sections d'action à $L_\alpha$ et $c$ fixés,
le quotient $\rho=\hbar_b/\hbar_a$ donne
$G_b=G_a/\rho$, $m_{P,b}=\rho m_{P,a}$ et
$E_{P,b}=\rho E_{P,a}$ ; $L_\alpha$ et $t_P$ restent identiques.
Cette comparaison entre sections maintient la distinction avec une
évolution temporelle des constantes.

\subsection{Covariance de l'archive et transformation des unités}

Soient \(T:V_{\rm in}\to V'_{\rm in}\) et
\(S:V_{\rm out}\to V'_{\rm out}\) unitaires. Le transport conjoint
des données est
\[
 B'=SBT^*,\quad P'=SPS^*,\quad P_i'=SP_iS^*,\quad X'=TXT^*.
\]
L'espérance transportée \(\mathcal E'(A')=S\mathcal E(S^*A'S)S^*\) fixe \(P_i'\) et est bimodulaire dans l'algèbre marquée transformée. Par substitution,
\begin{align*}
 \mathsf U'&=S\mathsf UT^*,&D'&=SDT^*,\\
 R'_{X'}&=SR_XS^*,&H'_{X'}&=SH_XS^*,&
 \Lambda'_{X'}&=S\Lambda_XS^*.
\end{align*}
La conjugaison conserve les égalités de R1--R8 et le même coefficient
\(G\). Les marques, orientations et registres sont transportés avec les
bases. Les permutations d'incidences constituent un cas particulier.

L'incorporation d'un espace auxiliaire de mémoire \(\mathcal F\)
transporte \(P\) vers \(P\otimes I_{\mathcal F}\),
\(\mathsf U\) vers \(\mathsf U\otimes I_{\mathcal F}\) et
\(\Delta_\Omega\) vers \(\Delta_\Omega\otimes\operatorname{id}\).
L'inscription de \(C_0C_0^*\otimes I_{\mathcal F}\) est
\(r_N I\otimes I_{\mathcal F}\) ; le coefficient conserve le cardinal
marqué original.

Pour préciser la covariance dimensionnelle, soient \(\ell\), \(\tau\) et \(m\) les facteurs positifs des nouvelles unités de longueur, de temps et de masse par rapport aux précédentes. Les coordonnées numériques satisfont
\[
 L'=L/\ell,\qquad c'=c\tau/\ell,\qquad
 \hbar'_{\rm ret}=\hbar_{\rm ret}\tau/(m\ell^2).
\]
Ainsi,
\[
 \frac{(c')^3(L')^2}{\hbar'_{\rm ret}}
 =G\frac{m\tau^2}{\ell^3},
\]
qui est la transformation d'une grandeur de dimension
\(\mathrm{longitud}^3\mathrm{masa}^{-1}\mathrm{tiempo}^{-2}\).
La section \(L_*\) se transforme comme \(L\), tandis que \(\alpha\)
et \(r_N\) restent sans dimension.

\subsection{Élimination exacte de la mémoire couplée}

Soient \(\mathcal K\) et \(\mathcal H\) des espaces de Hilbert finis,
\(W=W^*>0\) un poids global et
\(\mathcal C:\mathcal K\to\mathcal H\) surjective. L'action
des résidus \(r\), avec un défaut de frontière \(z\), est
\[
 \mathcal S(r)=\mathfrak a\,r^*Wr,
 \qquad \mathcal Cr=z,\qquad \mathfrak a>0.
\]
Le poids global conserve ses blocs croisés entre incidences. Définissons \(\mathcal R=\mathcal CW^{-1}\mathcal C^*\). Pour \(v\ne0\), l'injectivité de \(\mathcal C^*\) donne
\[
 v^*\mathcal Rv=\|W^{-1/2}\mathcal C^*v\|^2>0.
\]
Par conséquent, \(\mathcal R\) est inversible. Le résidu d'action minimale est
\[
 r_*(z)=W^{-1}\mathcal C^*\mathcal R^{-1}z.
\]
On vérifie \(\mathcal Cr_*=z\). Chaque résidu admissible s'écrit \(r=r_*+k\), avec \(\mathcal Ck=0\), et \(k^*Wr_*=k^*\mathcal C^*\mathcal R^{-1}z=0\). En conséquence,
\begin{equation}
 \mathcal S(r)=\mathfrak a\,z^*\mathcal R^{-1}z
             +\mathfrak a\,k^*Wk.
 \label{g26:eq:accion-residuo-es}
\end{equation}
La positivité de \(W\) prouve l'unicité du minimum. Le couple
\((z,k)\) conserve et reconstruit le résidu complet.

Pour une route avec les transports \(U_j\), les résidus
\(r_j=f_j-U_jf_{j-1}\) se reconstruisent par
\(f_j=U_jf_{j-1}+r_j\). L'application \(\mathcal C\) réunit les
transports ultérieurs de chaque résidu vers l'extrémité. Si ses
blocs sont \(\mathcal C_j\), alors
\[
 \mathcal R=\sum_{j,k}\mathcal C_j(W^{-1})_{jk}\mathcal C_k^*.
\]
Cette expression conserve les corrélations entre tronçons de la route. Pour un transport longitudinal inversible \(D\), la frontière \(y=Dz\) a pour réponse
\[
 \mathcal R_L=D\mathcal RD^*,\qquad
 \mathcal S_{{\rm ef},L}(y)
 =\mathfrak a\,y^*(D\mathcal RD^*)^{-1}y
 =\mathfrak a\,z^*\mathcal R^{-1}z.
\]
On obtient la même action en éliminant d'abord la mémoire ou en transportant
d'abord la frontière. Le préfacteur conserve la normalisation de l'action
source ; lorsque cette action porte \(\mathfrak a=\hbar_{\rm ret}\), la
réduction transmet la même section.

\subsection{Réduction conjointe du rayon, de l'énergie et de la longueur conjuguée}

Écrivons le même caractère positif dans une décomposition en frontière
et mémoire :
\[
 Y=\begin{pmatrix}A&F\\F^*&C\end{pmatrix},\qquad C>0,
 \qquad Y_{\rm ef}=A-FC^{-1}F^*.
\]
L’identité
\[
 \begin{pmatrix}b\\y\end{pmatrix}^{\!*}Y
 \begin{pmatrix}b\\y\end{pmatrix}
 =b^*Y_{\rm ef}b+\eta^*C\eta,
 \qquad \eta=y+C^{-1}F^*b,
\]
prouve \(Y_{\rm ef}>0\), caractérise l'extension d'énergie minimale et conserve la reconstruction \(y=\eta-C^{-1}F^*b\). Pour \(s>0\), la substitution dans le complément de Schur donne \(\operatorname{Schur}(sY)=sY_{\rm ef}\). En conséquence,
\[
 R_{\rm ef}=LY_{\rm ef},\qquad
 H_{\rm ef}=E_{L}Y_{\rm ef}.
\]
Pour calculer la longueur conjuguée comprimée, on résout
\(Y(x,y)=(f,0)\). La seconde équation donne
\(y=-C^{-1}F^*x\), et la première donne
\(x=Y_{\rm ef}^{-1}f\). Le bloc de frontière de \(Y^{-1}\) est
donc \(Y_{\rm ef}^{-1}\). Ainsi,
\begin{equation}
 \Lambda_{\rm b}=LY_{\rm ef}^{-1},\qquad
 R_{\rm ef}\Lambda_{\rm b}=L^2I,\qquad
 R_{\rm ef}=\frac{L}{E_{L}}H_{\rm ef}.
 \label{g26:eq:schur-reciproco-es}
\end{equation}
La réduction agit sur les mêmes blocs dans les deux lectures et
conserve le résidu \(\eta\). Le caractère effectif peut dépendre de
la mémoire sous forme matricielle et le coefficient \(L/E_{L}\) reste fixe.

\subsection{Mémoire dynamique et norme temporelle induite}

L'élimination dynamique conserve une dépendance spectrale supplémentaire. Dans une représentation stationnaire, soit
\[
 H=\begin{pmatrix}H_{\rm bb}&V\\V^*&H_{\rm ii}\end{pmatrix},
 \qquad H_{\rm ii}>0.
\]
Pour \(z\) dans le domaine résolvant pertinent, l'élimination de la
composante intérieure dans \((H-zI)(b,y)=(f,0)\) donne
\begin{equation}
 \mathcal K_H(z)=H_{\rm bb}-zI
       -V(H_{\rm ii}-zI)^{-1}V^*.
 \label{g26:eq:resolvente-memoria-es}
\end{equation}
Lorsque les deux inverses existent, le bloc de frontière de
\((H-zI)^{-1}\) est \(\mathcal K_H(z)^{-1}\). Pour
\(|z|<\lambda_{\min}(H_{\rm ii})\), la série convergente de la
résolvante fournit
\begin{align}
 \mathcal K_H(z)&=H_{\rm est}-zZ-z^2VH_{\rm ii}^{-3}V^*-\cdots,
 \label{g26:eq:expansion-resolvente-es}\\
 H_{\rm est}&=H_{\rm bb}-VH_{\rm ii}^{-1}V^*,\qquad
 Z=I+VH_{\rm ii}^{-2}V^*\geq I.
 \label{g26:eq:norma-z-es}
\end{align}
La positivité découle de \(Z-I=(H_{\rm ii}^{-1}V^*)^*(H_{\rm ii}^{-1}V^*)\). L'extension statique \((b,-H_{\rm ii}^{-1}V^*b)\) a pour carré de norme \(b^*Zb\) ; l'énergie et la norme temporelle proviennent de la même extension. Le paramètre physique est \(z=\hbar_{\rm ret}\omega\).

Au premier ordre temporel, \(T=Z^{1/2}\) et \(\psi=Tb\) transforment
le couple \((H_{\rm est},Z)\) en
\((T^{-*}H_{\rm est}T^{-1},I)\). La section
\(\hbar_{\rm ret}\) demeure fixée. Si \(T\) dépend du temps,
la transformation conserve également le terme \(\dot T b\) de
\(\dot\psi=\dot T b+T\dot b\).

La proportionnalité radiale est conservée dans la résolvante complète.
En écrivant \(\chi_{\rm rad}=L/E_{L}\) et \(R=\chi_{\rm rad} H\), la substitution de
chaque bloc démontre
\begin{equation}
 \mathcal K_R(\chi_{\rm rad} z)=\chi_{\rm rad}\mathcal K_H(z).
 \label{g26:eq:resolvente-proporcional-es}
\end{equation}
Cette identité transporte simultanément l'opérateur et le paramètre spectral, et conserve chaque ordre de la mémoire dynamique. Pour les formes effectives de \eqref{g26:eq:schur-reciproco-es}, la normalisation canonique exige le transport dual de la longueur inverse :
\begin{align*}
 H_{\rm can}&=T^{-*}H_{\rm ef}T^{-1},&
 R_{\rm can}&=T^{-*}R_{\rm ef}T^{-1},&
 \Lambda_{\rm can}&=T\Lambda_{\rm b}T^*.
\end{align*}
Multiplier ces expressions donne
\[
 R_{\rm can}=\frac{L}{E_{L}}H_{\rm can},\qquad
 R_{\rm can}\Lambda_{\rm can}=L^2I.
\]
La démonstration admet \(T\), \(R_{\rm ef}\) et
\(\Lambda_{\rm b}\) non commutatifs. La norme temporelle induite et
la longueur duale conservent conjointement le même \(G\).

\subsection{Raffinement, domaines spectraux et portée}

Soient \(J\) et \(K\) des inscriptions isométriques entre deux niveaux qui
satisfont \(X'J=JX\) et \(\mathsf U'J=K\mathsf U\),
avec la même section intensive \(L\). Alors
\[
 R'K=L\mathsf U'X'\mathsf U'^*K
     =KL\mathsf U X\mathsf U^*=KR,
\]
et le même calcul démontre \(H'K=KH\). L'entrelacement des
inverses se formule sur les domaines correspondants. Lorsque le système
compatible possède une limite autoadjointe positive et injective \(X\),
le calcul spectral définit \(Y=\mathsf U X\mathsf U^*\),
\(R=LY\), \(H=E_{L}Y\) et \(\Lambda=LY^{-1}\). La proportionnalité
\(R=(L/E_{L})H\) vaut sur \(\operatorname{Dom}(Y)\), tandis que
\(R\Lambda=L^2I\) et \(H\Lambda=\hbar_{\rm ret}cI\) valent sur
\(\operatorname{Dom}(Y^{-1})\), avec un prolongement borné de leurs
produits. Les coupures spectrales
\(\mathbf1_{[\varepsilon,M]}(Y)\) vérifient les identités avant
le passage à la limite dans ces domaines. Une longueur d'arête raffinée
\(L/9\) incorpore son propre facteur d'échelle ; la section intensive
se conserve selon l'entrelacement déclaré.

La classification distingue les transformations qui préservent la réalisation de celles qui modifient une règle constitutive. Les pondérations diagonales, les équivalences unitaires, l'archive auxiliaire et la réduction conjointe conservent \(G\). Le changement de résolution marquée modifie R1 ; le remplacement de \(r_N\) par \(\sqrt{r_N}\), de \(\alpha^{16}\) par son carré ou de la composante inscrite par la norme totale modifie R4. La remise à l'échelle isolée du rayon change R5, et celle de l'action change R6. R8 maintient la distinction entre rayon réduit, rayon de Schwarzschild et changement d'unités.

On obtient ainsi l'unicité mathématique de la réalisation R1--R8 : les données
générées et le caractère fixent la réponse, et la composition commune fixe
\(G\). Les règles longitudinale et métrique figurent explicitement dans
cette conclusion comme compositions incorporées ; l'identification
gravitationnelle conserve son statut physique. La preuve couvre la classe
déclarée avec ses définitions constitutives et ses transformations
compatibles. Les vérifications rationnelles du développement vérifient
des cas finis des identités exposées ; la démonstration générale
réside dans les équations et les arguments précédents.
% END CR28_RADIAL_FULL

